%------------------------------------------------------------------------------
% \PART 4
%------------------------------------------------------------------------------
\chapter[Разработка структуры IP-блока нейронной сети для распознавания рукописных цифр]
{РАЗРАБОТКА СТРУКТУРЫ IP-БЛОКА НЕЙРОННОЙ СЕТИ ДЛЯ РАСПОЗНАВАНИЯ РУКОПИСНЫХ ЦИФР}

%------------------------------------------------------------------------------
% \PART 4.1 Text
%------------------------------------------------------------------------------
\section{Описание структурной схемы}\par
\hspace*{12.5 mm}Принцип работы данной нейронной сети основан на комбинации 
обучаемого двумерного разделимого преобразования и полносвязной нейронной сети. 
Такая архитектура позволяет эффективно обрабатывать изображения, снижая 
вычислительные затраты при аппаратной реализации.

Для окончательной классификации используется полносвязный слой с функцией 
активации {softmax}. Данный слой преобразует выходные значения в вероятностное 
распределение по классам, позволяя определить, к какой цифре относится входное 
изображение.

Вычислительное ядро разработанного устройства состоит из 28 MAC 
(Multiply-ACcumulate) ядер, распределенных по специализированным блокам. Эти 
блоки организованы таким образом, чтобы обеспечить эффективное выполнение 
операций умножения входного вектора изображения на матрицу весов слоя
нейронной сети.

Десять MAC-ядер из общего массива используются на финальном тапе вычислений и 
непосредственно участвуют в формировании входных данных для функции активации 
softmax. Эти ядра получают входные данные из трех независимых блоков памяти, 
что позволяет оптимизировать процесс выборки весов и повысить скорость 
вычислений. Остальные 18 MAC-ядер распределены по блокам, где каждая 
вычислительная единица получает данные из двух запоминающих устройств.

IP-блок устройства предназначен для интеграции в систему на кристалле (СнК) и 
взаимодействует с процессорной системой через AXI-uP интерфейсы. СнК – это 
интегральная схема, которая объединяет в одном чипе все основные компоненты
вычислительной системы: процессор, ОЗУ (RAM) и ПЗУ (ROM), графический процессор
(GPU), контроллеры периферии. Общая структура представлена на 
рисунке~\ref{fig:str_all}.

\insertfigure{str_all}{pic/str.png}{Пример использования IP-блока нейронной сети}{16cm}

Последовательность обработки данных организуется следующим образом:

  1 Исходное изображение размером \(28\times28\) пикселей поступает на вход 
нейронной сети через переходник интерфейсов AXI-uP. Данные интерфейсы 
используется для взаимодействия с процессорной системой и обеспечивают передачу 
изображения в IP-блок. Пиксели изображения принимаются последовательно и 
записываются во внутреннюю память IP-блока, предназначенную для хранения 
изображения на протяжении всего цикла обработки.

  2 Каждый отдельный пиксель изображения поочередно передаётся на вход массива 
из 28 умножающих-накопительных ядер (MAC), которые одновременно выполняют 
операции умножения входного значения пикселя на соответствующий весовой 
коэффициент. Весовые коэффициенты извлекаются из специализированных блоков 
памяти весов, которые хранят параметры, полученные в результате 
предварительного обучения модели.

  3 За согласованность и последовательность операций отвечает управляющий 
модуль. Он обеспечивает адресацию ячеек памяти изображения и весов, организует 
подачу данных в вычислительные блоки и отслеживает фазы обработки. Управляющее 
устройство также координирует этапы переключения между слоями, а также 
формирует сигналы запуска и завершения вычислений.

  4 После завершения расчетов в первом скрытом слое (LST-1), результат – новое 
представление входного изображения — передаётся во второй уровень вычислений. 
Здесь используются 10 MAC-ядер, каждое из которых обрабатывает соответствующий 
вектор признаков, поступающий от первого слоя. В результате получается выходной 
вектор размерности 10, где каждый элемент отражает степень принадлежности 
изображения к соответствующему классу (от 0 до 9).

  5 Полученный выходной вектор подаётся в блок активации softmax. Этот модуль 
преобразует значения в вероятностную форму, нормируя их так, чтобы сумма 
элементов вектора была равна единице. На выходе формируется распределение 
вероятностей принадлежности изображения к каждому из 10 классов.

  6 Индекс элемента вектора softmax, имеющий наибольшее значение (то есть 
наиболее вероятный класс), определяется и передаётся через интерфейс AXI4-lite 
обратно в процессорную систему. Этот интерфейс реализует простую и эффективную 
регистровую модель доступа, обеспечивая передачу управляющих команд и данных 
между IP-блоком нейронной сети и процессором. Интерфейс uP (user processor 
interface) позволяет интегрировать модель в более широкую систему на кристалле, 
облегчая её взаимодействие с внешним программным обеспечением и обеспечивая 
гибкость конфигурации \cite{chu_fpga}.

Основные сигналы uP-интерфейса:

  1 {ADDR} – адрес регистра, к которому выполняется обращение.

  2 {DATA\_IN} – входные данные, записываемые в регистр.

  3 {DATA\_OUT} – выходные данные, считываемые из регистра.

  4 {RD} – сигнал чтения из регистра.

  5 {WR} – сигнал записи в регистр.

  6 {CLK} – синхросигнал интерфейса.

  7 {RESET} – глобальный сброс интерфейса.

  8 {READY} – флаг готовности к передаче данных.

  Временная диаграмма работы uP-интерфейса представлена на 
рисунке~\ref{fig:uP}.



  AXI4-Lite – это упрощённый вариант AXI4-интерфейса, предназначенный для 
управления и обмена небольшими порциями данных. В отличие от AXI4, он 
поддерживает только однослотовые транзакции без разделения на несколько 
пакетов\cite{vivado_axi}.

Основные сигналы AXI4-Lite:

    1 {AWADDR} – адрес записи.

    2 {AWVALID} – флаг валидности адреса записи.

    3 {WDATA} – записываемые данные.

    4 {WVALID} – сигнал валидности данных.

    5 {BREADY} – подтверждение завершения записи.

    6 {ARADDR} – адрес чтения.

    7 {ARVALID} – флаг валидности адреса чтения.

    8 {RREADY} – готовность принять данные.

    9 {RDATA} – считанные данные.

    10 {RVALID} – сигнал валидности данных.

    11 {ACLK} – системный тактовый сигнал.

    12 {ARESETN} – асинхронный сброс, активный по нулю.

\insertfigure{uP}{pic/up.png}{Временная диаграмма транзакций чтения и записи интерфейса uP}{16cm}

  Временная диаграмма работы AXI4-lite-интерфейса представлена на 
рисунке~\ref{fig:axi4}.
    
\insertfigure{axi4}{pic/axi.png}{Временная диаграмма транзакций чтения и записи интерфейса AXI4-lite}{16cm}

Оба интерфейса широко применяются в проектировании цифровых систем. 
uP-интерфейс удобен для работы с простыми IP-блоками, тогда как 
AXI4-Lite предоставляет стандартизированный метод взаимодействия с 
процессорными системами, такими как ARM-ядра на FPGA.\@ 

Электрическая структурная схема IP-блока нейронной сети прямого распространения 
приведена в графическом материале ГУИР 431289.002 Э1.

%------------------------------------------------------------------------------
% \PART 4.2 Text
%------------------------------------------------------------------------------
\section{Программная реализация эталонной модели нейронной сети на языке Python}
\hspace*{12.5 mm}Программная реализация нейронной сети выполняется с 
использованием языка Python.  
В качестве эталонной модели реализуется нейронная сеть на основе библиотеки 
{fixpoint}, позволяющая провести сравнение точности вычислений при 
использовании чисел с фиксированной запятой и чисел с плавающей запятой.

Эталонная модель представляет собой нейросеть с двумя скрытыми слоями и 
активационной функцией {tanh}.  

Для аппаратной реализации на FPGA и эталонной модели числа представляются в 
формате с фиксированной запятой (Fixed-Point). В отличие от чисел с плавающей 
запятой (Floating-Point), данный формат обеспечивает более эффективное 
использование аппаратных ресурсов, таких как блоки DSP и BRAM.\@

В данной работе используется фиксированное представление с разрядностью $Q6.7$,
где:

  – 6 бит отведены под целую часть числа.

  – 7 бит используются для дробной части.

Таким образом, числа в этом формате могут представлять значения в диапазоне 
$[-32, 31.996]$ с шагом $2^{-7} = 0.0078125$.

При переводе значений из плавающей запятой в фиксированную используется метод 
округления к ближайшему минимальному числу (округление вниз, или Round Down). 
Этот метод позволяет минимизировать ошибки округления и избежать 
систематического смещения, которое может возникать при округлении к ближайшему 
числу.

Такой метод округления позволяет повысить точность вычислений и уменьшить 
накопление ошибок при последовательных операциях, что критично для работы 
нейросети на FPGA.\@

Входные изображения размером \(28 \times 28\) проходят последовательную 
обработку:

  – Первый слой выполняет линейное преобразование входных данных с весами и 
смещением, после чего применяется функция активации тангенс.
    
  – Второй слой аналогично выполняет линейное преобразование, используя 
параметры весов и смещений, после чего применяется функция активации тангенс.
    
  – Выходной слой выполняет линейное преобразование входных данных с весами и 
смещением, а затем классифицирует изображение, вычисляя вероятности классов с 
помощью softmax.

Для проверки корректности работы алгоритма в числах с фиксированной запятой 
разрабатывается программная эмуляция аппаратного вычислителя. Эмуляция 
реализуется с использованием классов:

  1 {MAC} – блок умножения и аккумуляции с фиксированной точностью.
    
  2 {MEM} – память для хранения весов, смещений и промежуточных значений.
    
  3 {TANH} – реализация функции гиперболического тангенса в фиксированной 
точности.

Обработка входных данных в модели с фиксированной запятой включает следующие 
этапы:

  1 Загрузка изображения из базы данных MNIST в память.

  2 Преобразование входных данных согласно формату Q6.7.
    
  3 Выполнение первого преобразования и применение функции тангенс.

  4 Обработка данных во втором слое с аналогичной активацией.

  5 Вычисление выходных значений нейросети и применение softmax.

Сравнение выходных значений между моделью с фиксированной точностью и моделью с
плавающей точкой представлено на рисунке~\ref{fig:dif}.  

\insertfigure{dif}{pic/difpng.png}{Сравнение результатов вычислений}{16cm}

Графический анализ разности значений в обработке одного и того же изображения 
позволяет оценить влияние округлений и возможные потери точности.

Разработанная эталонная модель на Python позволяет проверить корректность 
вычислений перед аппаратной реализацией, а также провести тестирование 
точности чисел с фиксированной запятой.

Python описание нейронной сети с использованием библиотеки fixpoint приведено в
приложении В.

%------------------------------------------------------------------------------
% \PART 4.3 Text
%------------------------------------------------------------------------------
\section{Разработка операционной части нейронной сети}
\hspace*{12.5 mm}Нейронная сеть есть объединение управляющего и операционного
автоматов, что показано на рисунке~\ref{fig:operat}. Такое разбиение 
соответствует классическому подходу к построению цифровых вычислительных 
устройств, при котором логика работы системы делится на управление и выполнение.

\insertfigure{operat}{pic/operational.png}{Представление нейронной сети в виде управляющего и операционного автоматов}{16cm}

Элементарное действие, выполняемое в операционном автомате, называется 
микрооперацией. Пусть множество микроопераций, выполняемых под воздействием 
сигналов управляющего автомата, обозначается как \( Y = y_1, \dots, y_N\).  

Для выполнения микрооперации \( y_i \) (\( i = 1, \dots, N \)) необходимо 
появление соответствующего управляющего сигнала, принимающего значение единицы. 
Если в операционном автомате одновременно выполняется несколько микроопераций, 
они образуют микрокоманду.  

Множество микрокоманд в совокупности с логическими условиями 
\(X = x_1, \dots, x_N \), определяющими последовательность их выполнения, 
составляют микропрограмму. Логические условия соответствуют осведомительным 
сигналам, которые формируются в операционном автомате.

IP-блок нейронной сети должен быть спроектирован как независимый компонент 
системы, обеспечивающий модульность и возможность интеграции в различные 
вычислительные среды. Для этого необходимо, чтобы блок обладал четко 
определенным и простым интерфейсом, что позволит эффективно взаимодействовать 
с другими компонентами системы.

Интерфейс IP-блока должен обеспечивать передачу входных данных, управление 
процессом вычислений и выдачу результатов. Важно, чтобы структура 
взаимодействия соответствовала требованиям к гибкости и масштабируемости 
системы. 

На рисунке~\ref{fig:interface} представлен требуемый вид интерфейса 
IP-блока, демонстрирующий основные входные и выходные сигналы.

Входные сигналы включают:

– clk – тактовый сигнал, синхронизирующий работу IP-блока;

– rst\_n – активный по низкому уровню сигнал сброса, приводящий систему в исходное 
состояние;

– start\_i – управляющий сигнал запуска обработки входных данных;

– addr\_i – адресный сигнал, указывающий на текущую ячейку памяти изображения;

– data\_i – входные данные, представляющие собой значение пикселя изображения, 
поступающего в IP-блок.

Выходные сигналы включают:

– num\_o – выходной сигнал, содержащий распознанное значение цифры (от 0 до 9), 
представляющее итог работы нейронной сети;

– rdy\_o – сигнал готовности, свидетельствующий о завершении обработки и 
готовности результата к чтению со стороны процессорной системы.

\insertfigure{interface}{pic/interface.png}{Интерфейс IP-блока нейронной сети}{9cm}

%------------------------------------------------------------------------------
% \PART 4.4 Text
%------------------------------------------------------------------------------
\section{Проектирование функциональной схемы нейронной сети}
\hspace*{12.5 mm}При разработке IP-блока нейронной сети необходимо определить 
его функциональную структуру, обеспечивающую корректное выполнение вычислений, 
взаимодействие с внешними компонентами системы и эффективное управление 
процессом обработки данных.

Функциональная схема IP-блока должна описывать основные модули и их 
взаимодействие, включая блоки обработки входных данных, вычислительное ядро 
нейронной сети, управляющий модуль и интерфейсы для связи с внешней средой. 
Важно учитывать требования к производительности, ресурсоемкости и 
масштабируемости, чтобы обеспечить возможность интеграции IP-блока в различные 
аппаратные платформы.

% Также на функциональной схеме показаны микрооперации и логические условия. 
Данная нейронная сеть состоит из нескольких основных блоков. Существует два 
типа блоков обработки пикселя, которые отличаются количеством блоков памяти,
которые должны обрабатываться мультиплексором и поступать в качестве входных 
данных на умножитель.

После вычисления полносвязного слоя для конкретной строки или столбца данные 
поступают в регистр обработки, состоящий из 28 чисел разрядностью 13 бит. 
Затем данные из этого регистра последовательно поступают в регистр данных 
тангенса, из которого они отправляются в блок аппроксимации тангенса 
гиперболического.

Также используются два пяти разрядных и один десяти разрядный счетчики, которые
используются для генерации адреса в память весов и в память изображения.

Блок функции активации softmax реализован на принципе сравнения всех десяти 
входных классов и выбора наибольшего значения, что будет соответствует наиболее
вероятному значению. В результате на выход поступает индекс максимального 
элемента.

Для управления работой и генерирования микроопераций на основе логических 
условий используется микропрограммный автомат.

Электрическая функциональная схема IP-блока нейронной сети прямого 
распространения приведена в графическом материале ГУИР 431289.003 Э1. 

Электрическая функциональная схема блоков обработки пикселей приведена в 
графическом материале ГУИР 431289.004 Э1.

%------------------------------------------------------------------------------
% \PART 4.5 Text
%------------------------------------------------------------------------------
\section{Построение алгоритма работы нейронной сети}
\hspace*{12.5 mm}На данном этапе проектирования создается алгоритм, 
определяющая последовательность выполнения операций внутри IP-блока нейронной 
сети. Она описывает порядок обработки входных данных, передачу информации между 
модулями и выполнение вычислений.

Алгоритм строится на основе дискретных состояний, в которых выполняются 
конкретные микрооперации. Каждое состояние характеризуется активными 
управляющими сигналами, обеспечивающими выполнение определенной части 
вычислений.

Основные этапы работы алгоритма включают:

  1 Загрузку изображения в память – входные данные загружаются в буфер 
оперативной памяти, откуда они будут переданы в вычислительные блоки.

  2 Выполнение операций LST-преобразования – производится начальная обработка 
изображения, включающая нормализацию и предварительное выделение признаков.

  3 Передача данных в полносвязный слой – результаты преобразования передаются 
в полносвязный слой нейросети для дальнейшей обработки.

  4 Выполнение финальной классификации – производится анализ данных нейросетью
и определение класса входного изображения.

  5 Отправка данных в систему обработки – результаты классификации передаются
в систему для последующего использования.

Алгоритм работы нейронной сети по распознаванию рукописных цифр представлен в 
графическом материале ГУИР 431289.005 ПД. 

%------------------------------------------------------------------------------
% \PART 4.6 Text
%------------------------------------------------------------------------------
\section{Проектирование управляющей части}
\hspace*{12.5 mm}Проектирование управляющей части вычислительного устройства 
выполняется с применением метода структурирования конечных состояний, который 
широко используется при разработке цифровых систем управления.
Используется канонический метод структурного синтеза микропрограммного автомата 
(МПА), основанный на синхронной логике. Данный метод обеспечивает формальное 
построение управляющей логики устройства на основе конечного множества 
состояний, переходов между ними и соответствующих управляющих сигналов.

МПА делится на два основных структурных компонента:

1 Память управляющих микрокоманд, в которой хранится заданная последовательность 
управляющих слов. Каждое управляющее слово определяет активность тех или иных 
сигналов на каждом такте работы устройства, в зависимости от текущего 
состояния.

2 Комбинационная схема управления, которая отвечает за генерацию адреса 
следующей микрокоманды. Адрес определяется на основе текущего состояния 
автомата и внешних условий. Это позволяет гибко управлять переходами между 
состояниями в зависимости от динамики процесса обработки.

На рисунке~\ref{fig:cnn-pic} представлена обобщённая структурная схема 
такого микропрограммного автомата. 

\insertfigure{cnn-pic}{pic/cnt.png}{Представление управляющей части нейронной сети в виде композиции памяти и комбинационной схемы}{13cm}

Память автомата состоит из заранее выбранных элементов памяти, обеспечивающих 
хранение информации о текущем состоянии устройства. В качестве элементов памяти 
используются триггеры, которые позволяют надежно фиксировать и изменять 
состояние системы.  

Количество элементов памяти, необходимых для хранения состояний автомата, 
определяется по следующей формуле:
\begin{equation}
    N = \log_{2}{M},
\end{equation}

\noindentгде $M$ – число состояний микропрограммного автомата.  

Однако в практических реализациях часто применяется кодирование состояний с 
использованием кода Грея. Код Грея представляет собой двоичную систему 
кодирования, в которой любые два последовательных числа отличаются изменением 
только одного бита\cite{Tokheim2003}. Это позволяет уменьшить количество 
переключений триггеров при переходе между состояниями, снижая вероятность 
ошибок, возникающих из-за временных несовпадений сигналов.  

Дополнительно, в разработке автоматов управления применяется стратегия, при 
которой старший бит кодового слова указывает на нахождение автомата в состоянии
ожидания. Это состояние используется для инициализации системы, сброса 
счетчиков и минимизации нежелательных переходов. Применение старшего 
бита в качестве указателя на режим ожидания помогает уменьшить влияние гличей 
(glitch) на граф переходов. Гличи – это кратковременные нежелательные изменения
сигнала, возникающие при переключении состояний из-за различий во времени 
распространения сигналов через логические элементы~\cite{Katz1994}. Они могут 
приводить к некорректным переходам в автомате, вызывая ошибки в работе системы.  

В связи с этим, традиционная формула определения числа элементов памяти 
модифицируется следующим образом:
\begin{equation}
    N = \log_{2}{M} + 1 = \log_{2}{11} + 1 = 5,
\end{equation}

\noindentгде $M = 11$ – число состояний микропрограммного автомата.  

Каждое состояние автомата описывается представлением из 5 бит, 
что позволяет повысить надежность работы системы, уменьшить количество 
переключений триггеров и снизить вероятность ошибок, связанных с гличами.  

% Закодированные состояния автомата представлены в таблице 1.

В соответствии со списком микроопераций (Y) и логических условий (X), а также 
разработанной микропрограммой составим граф-схему алгоритма, которая 
представлена на рисунке~\ref{fig:gs}.

% \begin{table*}[ht]
%   Таблица 1 – Кодирование состояний автомата\par
%   \renewcommand{\arraystretch}{1.2}
%   \begin{tabularx}{\textwidth}{|c|>{\centering\arraybackslash}X|}
%     \hline
%     Состояние & Код Грея (\(T_4T_3T_2T_1T_0\)) \\ \hline
%     \(a_0   \)   & 00000                       \\ \hline
%     \(a_1   \)   & 10000                       \\ \hline
%     \(a_2   \)   & 10001                       \\ \hline
%     \(a_3   \)   & 10011                       \\ \hline
%     \(a_4   \)   & 10010                       \\ \hline
%     \(a_5   \)   & 10110                       \\ \hline
%     \(a_6   \)   & 10111                       \\ \hline
%     \(a_7   \)   & 10101                       \\ \hline
%     \(a_8   \)   & 11101                       \\ \hline
%     \(a_9   \)   & 11111                       \\ \hline
%     \(a_{10}\)   & 11110                       \\ \hline
%   \end{tabularx}
% \end{table*}

\insertfigure{gs}{pic/gs.png}{Граф-схема алгоритма}{16cm}

В соответствии с разработанной граф-схемой алгоритма составлена таблица 1, в 
которой представлены все переходы и их условия.

\begin{table*}[ht]
  Таблица 1 – Таблица переходов граф-схемы алгоритма\par
  \renewcommand{\arraystretch}{1.2}
  \begin{tabularx}{\textwidth}{|c|c|c|c|c|>{\centering\arraybackslash}X|}
    \hline
    h  & $a_m$    & $a_s$    & $X_h$               &                                           & \(Y_t\)                                \\ \hline
    1  & $a_0$    & $a_0$    & $\overline{x_0}$    &  –                                        & –                                      \\ \hline
    2  &          & $a_1$    & $x_0$               & $Y_1$                                     & $y_0, y_1, y_4, y_6, y_7, y_9, y_{10}$ \\ \hline
    3  & $a_1$    & $a_1$    & $\overline{x_1}$    &  –                                        &  –                                     \\ \hline
    4  &          & $a_2$    & $x_1$               &  –                                        &  –                                     \\ \hline
    5  & $a_2$    & $a_3$    &  1                  & $Y_2$                                     & $y_2$                                  \\ \hline
    6  & $a_3$    & $a_4$    &  1                  & $Y_3$                                     & $y_0, y_1, y_3$                        \\ \hline
    7  & $a_4$    & $a_4$    & $\overline{x_2}$    & $Y_4$                                     & $y_6$                                  \\ \hline
    8  &          & $a_5$    & $x_2$               &  –                                        &  –                                     \\ \hline
    9  & $a_5$    & $a_6$    &  1                  & $Y_5$                                     & $y_5$                                  \\ \hline
    10 & $a_6$    & $a_6$    & $x_4$ | $\overline{x_4}x_5$ | $\overline{x_4}\overline{x_5}x_6$ & – | $Y_6$ | $Y_7$  & $y_6, y_7$        \\ \hline
    11 &          & $a_7$    & ${x_2}$             & $Y_8$                                     & $y_8$                                  \\ \hline
    12 & $a_7$    & $a_3$    & $\overline{x_3}$    & $Y_9$                                     & $y_4$                                  \\ \hline
    13 &          & $a_8$    & $x_3$               & $Y_{10}$                                  & $y_{10}$                               \\ \hline
    14 & $a_8$    & $a_9$    &  1                  &  –                                        &   –                                    \\ \hline
    15 & $a_9$    & $a_9$    & $\overline{x_7}$    & $Y_{11}$                                  & $y_{11}$                               \\ \hline
    16 &          & $a_{10}$ & $x_7$               &  –                                        &   –                                    \\ \hline
    17 & $a_{10}$ & $a_0$    &  1                  &  –                                        &   –                                    \\ \hline
  \end{tabularx}
\end{table*}

%------------------------------------------------------------------------------
% \PART 4.7 Text
%------------------------------------------------------------------------------
\section{Построение графа переходов управляющего автомата}

\hspace*{12.5 mm}Граф переходов управляющего автомата позволяет представить 
граф-схему алгоритма в соответствии с условными обозначениями. Такая форма 
представления облегчает анализ работы системы и позволяет четко проследить 
взаимосвязь между различными этапами обработки данных. Граф переходов 
управляющего автомата представлен на рисунке~\ref{fig:gsa}.

Для каждого состояния определено мнемоническое обозначение, которое 
обеспечивает удобство представления и понимания процесса работы автомата. Все 
состояния кодируются в виде пятибитных значений, сформированных в коде Грея.

Состояния граф-схемы алгоритма объеденины в блоки:

  1 {INIT} – инициализация системы, установка начальных параметров, необходимых 
для корректной работы всех последующих этапов.

  2 {LOAD} – загрузка изображения.

  3 {LD\_MEM} – подготовка данных.

  4 {INIT\_CALC} – инициализация смещения.

  5 {CALC} – обработка строки/столбца.

  6 {RDY} – готовность результата обработки.

  7 {TG\_CALC} – вычисление тангенса.

  8 {UPD\_RCO} – смена слоя.

  9 {INIT\_OUT} – инициализация выходных смещений.

  10 {OUT\_CALC} – обработка выходного полносвязного слоя.

  11 {RDY\_OUT} – готовность детектирования.

\insertfigure{gsa}{pic/gsa.png}{Граф переходов}{11.5cm}
